\chapter{Introduction}
\label{ch:introduction}

\section{Motivation}
Intelligent Transportation Systems (ITS) heavily rely on accurate short-term and long-term traffic flow prediction to optimize routing, alleviate congestion, and reduce carbon emissions. Historically, traffic prediction models have relied on centralized architectures, where raw data collected from thousands of ubiquitous sensors is continuously transmitted to a central cloud server for training. While centralized models—ranging from traditional statistical methods like ARIMA to advanced Spatio-Temporal Graph Neural Networks (STGCN) and Sequence-to-Sequence (Seq2Seq) models—can achieve high accuracy, this paradigm poses two critical challenges: significant communication overhead and severe data privacy risks.

Federated Learning (FL) has emerged as a promising paradigm to address these challenges by allowing edge devices (e.g., roadside units or vehicular networks) to collaboratively train a shared model without exchanging their raw local data. Instead of data, only model parameters or gradients are communicated to the aggregation server. However, applying standard federated aggregation algorithms, such as FedAvg \cite{mcmahan2017}, to traffic forecasting introduces new obstacles. Traffic data is highly non-Independent and Identically Distributed (non-IID); sensors on a highway experience drastically different temporal patterns compared to those on residential arterials. When a single global model attempts to fit all these divergent patterns simultaneously, it often suffers from slow convergence and suboptimal accuracy. Furthermore, in realistic vehicular networks, client availability is dynamic, and the communication bandwidth required to frequently upload large deep learning models like Sequence-to-Sequence (Seq2Seq) GRUs remains prohibitively high \cite{zhu2025cmba}. 

To bridge this gap, there is a pressing need for a structured federated learning framework that can handle the extreme spatial-temporal heterogeneity of traffic data while aggressively mitigating communication costs and adapting to unreliable edge network conditions.

\section{Aim and Research Questions}
The primary aim of this thesis is to design, implement, and evaluate a novel \textbf{Hierarchical Federated Learning Framework} specifically optimized for communication-efficient traffic flow prediction in Intelligent Transportation Systems.

This work is guided by the following central research question:

\begin{quote}
\textit{Can clustering traffic sensors by temporal pattern similarity improve the accuracy and communication efficiency of federated learning for traffic flow prediction?}
\end{quote}

\noindent To systematically address this question, the thesis investigates four sub-questions:

\begin{enumerate}[label=\textbf{SQ\arabic*:}]
    \item How does DTW-based clustering compare to global federated averaging for heterogeneous traffic data?
    \item Does hierarchical knowledge sharing between clusters improve prediction accuracy over isolated cluster training?
    \item Can adaptive client selection reduce communication costs without significantly degrading model performance?
    \item Is traffic pattern similarity (DTW) a more effective basis for inter-cluster aggregation than geographic proximity?
\end{enumerate}

\section{Objectives}
To answer the above research questions, the thesis defines the following specific objectives:
\begin{enumerate}
    \item \textbf{Traffic-Pattern Aware Clustering:} To implement Dynamic Time Warping (DTW) to cluster traffic sensors based on temporal behavioral similarity rather than mere geographic proximity, ensuring that models specialize in homogeneous traffic patterns \cite{feng2024}.
    \item \textbf{Hierarchical Aggregation Strategy:} To develop a two-level aggregation scheme where intra-cluster models are trained locally (Level 1), and inter-cluster knowledge is shared periodically through a distance-weighted global average (Level 2), thereby balancing model specialization with spatial generalization.
    \item \textbf{Communication Efficiency and Robustness:} To reduce communication overhead through an adaptive, reliability-based client selection mechanism that dynamically selects participating nodes each round, coupled with quality-weighted aggregation that prioritizes higher-performing clients during model merging \cite{dai2024}.
    \item \textbf{Comprehensive Ablation and Benchmarking:} To rigorously validate the proposed framework on the real-world METR-LA dataset using a Sequence-to-Sequence GRU architecture \cite{zhu2025cmba}, directly comparing the hierarchical clustered approach against global FedAvg and centralized baselines.
\end{enumerate}

\section{Thesis Outline}
The remainder of this thesis is organized as follows:

\begin{itemize}
    \item \textbf{Chapter 2: Background and Literature Review.} This chapter provides the theoretical foundations of sequence modeling in traffic prediction, introduces the principles of Federated Learning, and critically reviews recent advancements in Spatio-Temporal FL, client selection, and communication sparsification techniques within Intelligent Transportation Systems.
    
    \item \textbf{Chapter 3: Methodology.} This chapter details the proposed Hierarchical Federated Learning framework. It provides an in-depth mathematical and architectural description of the DTW-based clustering algorithm, the Seq2Seq GRU local models, the adaptive client selection policy, the quality-weighted aggregation, and the two-tier distance-weighted hierarchical aggregation strategy.
    
    \item \textbf{Chapter 4: Ablation Study.} This chapter isolates the individual components of the proposed pipeline. It systematically evaluates the impact of clustering versus a flat global network, the contribution of hierarchical inter-cluster aggregation, the effect of adaptive client selection on communication costs, and the comparison between DTW-weighted and geography-weighted hierarchical sharing.
    
    \item \textbf{Chapter 5: Results and Discussion.} This chapter presents the final comparative performance of the complete framework against baseline models on the METR-LA dataset. It analyzes the results across multiple prediction horizons (15, 30, and 60 minutes), discusses the per-cluster performance variations, and evaluates the overall communication cost reductions achieved.
    
    \item \textbf{Chapter 6: Conclusion.} The final chapter summarizes the key findings and contributions of the research, acknowledges the limitations of the current implementation, and outlines promising directions for future work, including the integration of external topological metadata.
\end{itemize}
